\section{Packing: a smooth energy and its exact gradient}\label{sec:packing}
\stagebox{\statusrework}{\src{core/src/energy.rs}, \src{core/src/pack.rs}}

The sheets laid out by \cref{sec:sheets} fit their patches exactly. Packing
makes every petal \emph{want} to be larger than the room it has, then lets all
petals settle together by minimising one energy. Petals that cannot grow
outward must buckle and fold around each other, and that folding is the 3D
packing.

\subsection{Unknowns}
The unknowns are the positions $\x_i\in\R^3$ of all control points of all
petals, $X = (\x_1,\dots,\x_n)$. Point $i$ belongs to petal $p(i)$ on layer
$k(i)$ and has a fixed thickness $t_i$. The base row ($v=0$) of every petal is
pinned to the receptacle and never moves.

\subsection{The energy}
\[
  E(X) = E_{\mathrm{spring}} + E_{\mathrm{band}} + E_{\mathrm{env}} + E_{\mathrm{contact}}.
\]
Every term is a spring or a squared hinge $\tfrac12 k\,[\,\cdot\,]_+^2$, so $E$ is
continuously differentiable ($C^1$) and has an exact gradient.

\paragraph{Springs (stretch, shear, bending).} Each petal's control grid is
joined by springs $(a,b)$ with stiffness $k_{ab}$ and rest length $\ell_{ab}$:
\[
  E_{\mathrm{spring}} = \sum_{(a,b)} \tfrac12 k_{ab}\bigl(\lVert\x_a - \x_b\rVert - \ell_{ab}\bigr)^2,
  \qquad
  \frac{\partial}{\partial \x_a} = k_{ab}\,\frac{\lVert\x_a - \x_b\rVert - \ell_{ab}}{\lVert\x_a - \x_b\rVert}\,(\x_a - \x_b).
\]
\begin{center}\small
\begin{tabular}{@{}lll@{}}
\toprule
Spring & Joins & Stiffness \\
\midrule
stretch & neighbours along $u$ and along $v$ & 40 \\
shear & both diagonals of each grid cell & 10 \\
bending & points two apart along $u$ and along $v$ & $6\,(\bar t/\tmax)^3$ \\
\bottomrule
\end{tabular}
\end{center}
Here $\bar t$ is the petal's mean thickness. Bending stiffness of a plate grows
with the cube of its thickness, so thin petals crumple more easily than thick
ones. (Today $\bar t$ is one value per petal, not per point.)

Every rest length is the spring's laid-out length times the growth factor $g$:
\[
  \ell_{ab} = g\,\lVert\x^{0}_a - \x^{0}_b\rVert, \qquad g\in[1.12, 1.35]\text{ for peony}.
\]
Each petal therefore wants to be $g$ times larger in every direction, about
$g^2$ times its patch area, which it cannot achieve without leaving its slot or
pushing into its neighbours. This is what forces buckling.

\paragraph{Layer slot.} Each point should stay within its layer's slot, less its
own half-thickness. With half-width $h_i = \max\bigl(0, (w - t_i)/2\bigr)$ and
excess $\epsilon_i = \lvert\rt(\x_i) - a_{k(i)}\rvert - h_i$,
\[
  E_{\mathrm{band}} = \sum_i \tfrac12 k_{\mathrm{band}}\,[\epsilon_i]_+^2,
  \qquad
  \frac{\partial}{\partial \x_i} = k_{\mathrm{band}}\,[\epsilon_i]_+\,\operatorname{sign}\bigl(\rt - a_{k(i)}\bigr)\,\nabla\rt(\x_i).
\]

\paragraph{Envelope and core.} Points stay inside the sepals and outside the
stamens, each allowing for half its thickness:
\[
  E_{\mathrm{env}} = \sum_i \tfrac12 k_{\mathrm{env}}\Bigl(\bigl[\rt(\x_i) - (R_{\mathrm{env}} - \tfrac{t_i}{2})\bigr]_+^2 + \bigl[(R_0 + \tfrac{t_i}{2}) - \rt(\x_i)\bigr]_+^2\Bigr).
\]

\paragraph{Contact.} Two petals must keep their mid-surfaces at least their
combined half-thickness apart. For a point $\x_i$ and a triangle
$T = (a,b,c)$ of another petal's control mesh (each grid cell is split into two
triangles), let $\vect q = \lambda_a\x_a + \lambda_b\x_b + \lambda_c\x_c$ be the
closest point of $T$ to $\x_i$ (barycentric weights $\lambda$, computed by
Ericson's region test), $d = \lVert\x_i - \vect q\rVert$, and
\[
  d_{\min} = \tfrac12\Bigl(t_i + \tfrac13(t_a + t_b + t_c)\Bigr).
\]
Then
\[
  E_{\mathrm{contact}} = \sum_{(i,T)} \tfrac12 k_{\mathrm{contact}}\,[d_{\min} - d]_+^2 .
\]
Using point-to-triangle distance rather than point-to-point means a point cannot
slip through the gap between two control points of a neighbour. Every point is
tested against triangles of every other petal, so each pair of petals is
checked in both directions.

The gradient follows from the envelope theorem: because $\vect q$ minimises the
distance over $T$, moving $\vect q$ within the triangle does not change $d$ to
first order, so the weights $\lambda$ can be held fixed. With
$\vect n = \x_i - \vect q$ and $\kappa = k_{\mathrm{contact}}[d_{\min}-d]_+/d$,
\[
  \frac{\partial E}{\partial \x_i} = -\kappa\,\vect n,
  \qquad
  \frac{\partial E}{\partial \x_m} = \kappa\,\lambda_m\,\vect n \quad (m\in\{a,b,c\}).
\]
Using the triangle's mean thickness, instead of the thickness at $\vect q$,
keeps $d_{\min}$ independent of $X$, so no extra terms appear. The gradient is
exact wherever the closest point is unique.

\begin{center}\small
\begin{tabular}{@{}lc@{}}
\toprule
Constant & Value \\
\midrule
$k_{\mathrm{band}}$ & 400 \\
$k_{\mathrm{env}}$ (envelope and core) & 400 \\
$k_{\mathrm{contact}}$ & 1500 \\
\bottomrule
\end{tabular}
\end{center}

\paragraph{Checked by test.} A test squeezes a bud so that the slot, envelope
and contact terms are all active and compares the analytic gradient with
central finite differences ($h = 10^{-6}$) on 60 coordinates. It caught one
real bug during development.

\subsection{Breaking symmetry}
A perfectly laid-out sheet under uniform growth could stay balanced on its
shell without choosing a direction to buckle. Before relaxing, every unpinned
point is nudged along the shell normal $\nabla\rt/\lVert\nabla\rt\rVert$ by
\[
  \delta_i = 0.1\,w\,\frac{v_i}{\NV - 1}\,\xi_i, \qquad \xi_i\sim\Norm(0,1)\text{ from stream 2},
\]
growing from zero at the base to $0.1w$ at the tip. Because the nudge comes from
the plant's stream, all buds on a plant buckle the same way.

\subsection{Relaxation by FIRE}
The energy is minimised by FIRE (Fast Inertial Relaxation Engine), gradient
descent with inertia that is reset whenever it starts climbing. With
velocities $\vect v_i$ (pinned points have zero gradient and never move), each
iteration:
\begin{enumerate}[leftmargin=*]
\item computes $\nabla E$ and the power $P = -\nabla E\cdot \vect v$;
\item if $P > 0$, steers the velocity toward the downhill direction,
  $\vect v \leftarrow (1-\alpha)\vect v - \alpha\,\lVert\vect v\rVert\,\nabla E/\lVert\nabla E\rVert$,
  and after more than 5 such steps in a row sets
  $\Delta t \leftarrow \min(1.1\,\Delta t, \Delta t_{\max})$ and $\alpha \leftarrow 0.99\,\alpha$;
\item if $P \le 0$, stops: $\vect v \leftarrow 0$, $\Delta t \leftarrow \Delta t/2$, $\alpha\leftarrow 0.1$;
\item integrates $\vect v_i \leftarrow \vect v_i - \Delta t\,\nabla_i E$ and
  $\x_i \leftarrow \x_i + \Delta t\,\vect v_i$, with each point's step clipped to
  at most $\tmax/4$ so no petal can jump through another.
\end{enumerate}
Defaults: 2000 iterations, $\Delta t_{\max} = 0.03$, starting from
$\Delta t = \Delta t_{\max}/2$ and $\alpha = 0.1$.

\subsection{Finding contacts cheaply}
Testing every point against every triangle would be quadratic in the number of
petals. Instead a list of candidate (point, triangle) pairs closer than
\[
  \text{reach} = \tmax + \text{skin}, \qquad \text{skin} = 2\tmax,
\]
is built with a spatial hash, and reused until some point has moved more than
half the skin since the list was built. Contacts only matter below
$d_{\min}\le\tmax$, and a point and a triangle that each moved less than half
the skin cannot have closed more than one skin, so no contact is missed
between rebuilds. The hash is used only for lookup; the list is sorted
before use, so results never depend on hash order (\cref{sec:determinism}).

\subsection{What the packer reports}
After relaxation the core reports the energy terms before and after, the number
of contact-list rebuilds, the worst overlap between two petals
$\max(d_{\min} - d)/\tmax$, and the worst excursion outside a slot
$\max\epsilon_i/\tmax$. These are written into the flower JSON's \texttt{meta}
block.

\subsection{Current limits and the rework}\label{sec:packing-limits}
The current bud has two visible weaknesses: the buckling is gentle, so the bud
reads as neat nested layers rather than crumpled tissue, and the inner petals
stop partway up, leaving an empty dome at the top. The slot term
$E_{\mathrm{band}}$ is a large part of the first: it holds every point inside
its own layer.

Real peony buds swell like marshmallows and have very complex packing, and the
bud stage must be able to model that. The rework under way changes four
things:
\begin{itemize}[leftmargin=*]
\item \textbf{growth driven by the automaton}: per-patch random growth rates, so
  unequal growth forces crumples and ruffles (\cref{sec:ca});
\item \textbf{swelling}: the sepal envelope gives way when petals push on it
  (\cref{sec:shells});
\item \textbf{mixed layers}: petals may fold over each other and between
  layers, held apart only by their thickness, which loosens or removes the slot
  term;
\item \textbf{finer petals}: a much finer control grid per petal
  (\cref{sec:sheets}).
\end{itemize}
This section will be rewritten when the rework lands.
